% Obacht: Hier bitte nichts ändern, das verändert die Vorlage! Eure Eintragungen sollen ausschließlich in den Textfeldern erfolgen, die mit den runden Klammern versehen sind. Dafür könnt ihr entweder hier im Skript nach unten scrollen, starten würde dies ab Zeile 63. Ihr könnt aber auch die Pfeile nutzen, diese führen euch automatisch in die richtigen Zeilen des Codes, wo ihr die Daten in den runden Klammern dann mit euren Daten ersetzen könnt. Die Seitenzahlen im Inhaltsverzeichnis könnt ihr anklicken, dann werdet ihr automatisch an den entsprechenden Ort im Bericht weitergeleitet. Bei Fragen oder Anmerkungen meldet euch gerne per Discord oder Mail an Laura (laura-alia.kraski@student.uni-siegen.de). :)

\documentclass[11pt,a4paper,titlepage]{article}
\usepackage[left=2.0cm,right=2.5cm,top=2.5cm,bottom=2.8cm]{geometry}
\setlength{\parskip}{6pt}

\usepackage[utf8]{inputenc}
\usepackage[T1]{fontenc}
\usepackage[nottoc]{tocbibind}
\usepackage{amsmath}
\usepackage{dcolumn}
\usepackage{longtable}
\usepackage[font=bf]{caption}
\usepackage{amsfonts}
\usepackage{amssymb}
\usepackage{enumitem}
\usepackage{lscape}
\usepackage{setspace}
\usepackage{multirow}
\usepackage{graphicx}
\usepackage{microtype}
\usepackage{fancyhdr}
\newcommand{\changefont}[3]{\fontfamily{#1} \fontseries{#2} \fontshape{#3} \selectfont}
\usepackage[toc,page]{appendix}
\usepackage[linktocpage=true, colorlinks = true, linkcolor = blue, urlcolor  = blue, citecolor = blue, anchorcolor = blue]{hyperref}
% Sprachpaket ändern
\usepackage[german]{babel}

% Schriftarten ändern
\usepackage{textcomp}
%\usepackage{bookman} %bookman schrift
%\usepackage{mathpazo} %Palatino schrift
\usepackage{new cent} %Calibri

% Bibliography
\usepackage[style=authoryear-icomp, maxbibnames=9, maxcitenames=2, uniquelist=false, uniquename=false, isbn=false, doi=false, url=false, natbib=true, backend=biber]{biblatex}
\DeclareLanguageMapping{british}{british-apa}
\usepackage[autostyle]{csquotes}
\addbibresource{samplebib.bib} 

% Author information
\author{Autor*in}

\begin{document}

% Titelseite
\begin{titlepage}

\begin{minipage}[t]{\textwidth}
\includegraphics[width=\textwidth]{Picsart_23-05-11_12-59-33-529.jpg}
\end{minipage}


    \begin{center}
          
        \vspace*{1cm}
        
        \huge
        \textbf{Praktikumsbericht SkillsLab} \\
        
        \vspace{1.3cm}
        
        \textbf{Autor*in: (Deinen Namen hier einfügen)} \\
          \vspace{1.5cm}
        \normalsize
        Abgabe: (Datum hier eingeben)
     \vspace{3.5cm}

        Leistung im Modul: "Forschungsgrundpraktikum (5DBHSBA04)
       \vspace{0.8cm}
   
       Universität Siegen\\
       Bachelorstudiengang DBHS\\
       Fakultät V\\
      \vspace{1.7cm}
      
    \begin{tabular}{p{8cm}p{8cm}}
    Verantwortlicher: & Student*in: \\ 

    &  \\ 

    Christian Gießer & (Deinen Namen und deine Matrikelnummer hier einfügen)  \\ 

     E-Mail: christian.giesser@uni-siegen.de &  E-Mail: (Deine Mail-Adresse hier eintragen)\\ 
     \end{tabular}        
    \end{center}
\end{titlepage}

\setstretch{1.0} % Line spacing of 1.0

% Inhaltsverzeichnis
\tableofcontents
\addtocontents{toc}{~~\hfill\textbf{Page}\par}

\setstretch{1.0} % Line spacing of 1.0
\newpage
\section{Einführung} 
\label{sec:introduction}

(Hier sollst du das SkillsLab einmal vorstellen. Beschrieben werden soll hier, was das SkillsLab genau ist (Idee, Ziele und die Durchführung),  und deine Aufgabe während dem Forschungsgrundpraktikum sollte kurz erläutert werden.)
\newpage
\subsection{Aufgaben während dem Praktikum}
\label{sec:litrev}
(Hier sollst du deine Aufgabe genauer erläutern (Tutoren oder Sonderpositionen, hier ist noch keine genaue Erläuterung nötig. Es geht hierbei nur um eine kurze Einführung.)

\subsubsection{Vorbereitung}
(Hier sollst du umschreiben, wie du dich auf deine Aufgabe vorbereitet hast; Gab es Dinge, die du für das Forschungsgrundpraktikum noch lernen musstest?)
\newpage
\section{Station}
(Hier sollst du beschreiben, welche Station du betreut hast, und worum es bei der Station genau geht.)
\subsection{Ziele}
(Hier sollst du nochmal genauer erläutern, was das eigentliche Ziel der Station ist. Was soll den Studierenden hier näher gebracht werden? Was ist der Hintergrund?)
\subsection{Durchführung und Ablauf}
(Hier geht es um den Modellaufbau, die Durchführung. Der Ablauf soll hier einmal genauer geschildert werden.)
\subsection{Schwierigkeiten}
(Hier sollst du, vorausgesetzt es sind welche aufgetreten, die Schwierigkeiten erläutern. Gab es Dinge, die dir oder den Student*innen schwer gefallen sind? Wenn ja, warum? Was kann man in Zukunft besser machen?

\newpage
\section{Auswertung}
\subsection{Fragestellungen der Umfrage}
(Hier sollt ihr kurz nennen, welche Fragen es bei eurem Umfragebogen gab.)
\subsection{Ergebnisse der Umfrage}
(Hier sollt ihr die Ergebnisse aus der Umfrage vorstellen und Schlüsse aus dieser ziehen.)
\newpage
\section{Fazit}
(Hier kannst du deine Zeit während des Praktikums generell reflektieren. Was hat dir besonders gefallen, was hast du gelernt? Gibt es Dinge, die du dir gewünscht hättest oder die du verbessern würdest? 

\newpage
\section{Anhang}
(Hier sollt ihr Bilder von eurer Station einfügen. Ihr könnt den Modellaufbau und gegebenenfalls die Durchführung fotografieren und in euren Bericht packen. Genutzt wird dafür der Befehl \ includegraphics - bei Fragen meldet euch gern.)
\label{sec:models}


% Literaturverzeichnis
\newpage
\printbibliography[heading=bibintoc]




\end{document} 